\chapter{从大纲到全量构建}
\label{ch:build}

本章登记依赖顺序、各工作包的完成判据，以及主线在包与包之间如何衔接。没有日志的数据图不先画曲线。没有断言的页面不先发布为结果。文献边界没写明时，不把仿真残差写成物理定律。

\section{依赖}

\paragraph{要判定的量}
哪一步必须先于哪一步。并行只允许发生在互不作为输入的包之间。

\begin{center}
\small
\begin{tikzpicture}[
  node distance=5mm and 7mm,
  box/.style={draw, rounded corners, align=center, font=\small, inner sep=3pt, minimum width=24mm},
  arr/.style={-{Stealth}}
]
\node[box] (a) {符号与\\方程冻结};
\node[box, right=of a] (b) {六条命题\\的证明};
\node[box, right=of b] (c) {两节点\\与工程核};
\node[box, right=of c] (d) {无界面\\断言};
\node[box, below=8mm of a] (e) {多孔介质\\真值};
\node[box, right=of e] (f) {四策略\\日志};
\node[box, right=of f] (g) {D 图与\\交叉表};
\node[box, below=8mm of e] (h) {外部数据\\只作先验};
\node[box, right=of h] (i) {三维与\\机制视图};
\node[box, right=of i] (j) {PDF 与\\Word};
\draw[arr] (a)--(b);
\draw[arr] (b)--(c);
\draw[arr] (c)--(d);
\draw[arr] (d.south) |- (i.north);
\draw[arr] (a.south)--(e.north);
\draw[arr] (e)--(f);
\draw[arr] (f)--(g);
\draw[arr] (g.south) |- (j.north);
\draw[arr] (h)--(i);
\draw[arr] (i)--(j);
\end{tikzpicture}
\end{center}

上图里已经走完的是：符号与两节点工作方程、六条命题在其断言常数下的核对、十六条控制体的工程核、S1 日志、五张 D 图、清洁等温线作为淀粉先验的登记，以及由 TeX 源导出的 PDF 与 Word。尚未走完的是多孔介质真值积分器。因此从“多孔介质真值”指向“四策略日志”的箭，在本次只走到工程核，不到达带压力和蒸汽分数的场。视图可以与尚未完成的多孔介质并行开发，它不能先于断言发布，也不能用着色器里的颜色代替冷点。

Word 在 PDF 的公式编号稳定之后同步。不并行改两套正文。

\paragraph{本节不能推出的结论}
依赖图上的箭头画出来，不能写成箭头终点的计算已经运行。终点是否运行，以各包的完成判据为准。

\section{工作包}

\paragraph{要判定的量}
每个包的完成判据是哪一条可核对的记录，以及该包明确不做什么。

\begin{table}[ht]
\centering
\caption{工作包与本次记录}
\begin{tabularx}{\textwidth}{clXX}
\toprule
序 & 包 & 完成判据 & 本次记录 \\
\midrule
1 & 符号与单位 & 同一符号只指一个量 & 附录符号表已列出；全场与决策状态分开定义 \\
2 & 命题证明 & 六条命题有证明，Q1--Q6 只含证明几何 & 证明在第 \ref{ch:proof} 章；断言四条通过 \\
3 & 数值核 & 固定种子下轨迹可复现 & 正弦散列，无随机数发生器；种子 7 可复读 \\
4 & 断言 & 四条命题有独立测试 & 稳态极差 0；脉冲用 \(+1\) 准则；上界 1.143 小于 1.2；屏障不读后验 \\
5 & 多孔真值 & 含 GAB、各向异性、脂肪层、收缩的场 & 未运行。工程核不是该判据 \\
6 & 试验 & B0--B3 的日志与种子可复现 & S1 已运行，4860 组合；S0、S2、S3 未运行 \\
7 & 图与正文 & 有日志的 D 图画点，其余保留坐标 & D1、D6、D7、D10、D25 已画；其余 D 位无曲线 \\
8 & 演示 & 切面与导出和核的列一致 & 页面读取状态；浏览器副本不进主表 \\
9 & 双文件 & PDF 与 Word 编号一致 & 由同一 TeX 源导出 \\
\bottomrule
\end{tabularx}
\end{table}

包 5 的等温线先验使用第 \ref{ch:data} 章的清洁子集。肉类导热仍是文献区间，并在参数文件中与等温线来源分开。S1 的完成判据是四套策略的清单与种子可复现，不是轨迹条数本身。

主线上的检查点按下面的规则执行，不按“条数已经很多”执行。

\begin{itemize}
\item S0 的网格减半没有单独运行，因此不声明格式误差已低于某阈值。15 mm 上显式步长 \(2\,\mathrm{s}\) 距离扩散稳定性上限约 \(2.2\,\mathrm{s}\)，这一条写在离散格式里，不写成 S0 已经通过。
\item 命题断言不通过时，不把三维页面上的颜色当成冷点证据。本次四条断言在其自身常数下通过。
\item T2 尚未运行。不存在“T2 与 T1 排序相反”的日志，因此也不把两者平均成一个排序。
\item 演示若不能在时间轴上指出固定探针窗、在状态里指出热阻方差、在加热指令上指出削减是否动作，它与正文的闭环就不一致。本次加热器平衡在 \(196.5^\circ\mathrm{C}\)，削减阈值 \(230^\circ\mathrm{C}\) 不动作，页面不得把削减画成已经发生。
\end{itemize}

\paragraph{本节不能推出的结论}
九个包里已经有记录的那些，不能替尚未有记录的包出具结果句。包 5 未完成时，主结论的主语保持为工程核。

\section{成稿时的逻辑骨架}

全文只沿一条链写。章内可以展开推导，章间不另起一条互不引用的主题。

\[
\begin{aligned}
&\text{信息赤字}
\to\text{多尺度场}
\to\text{决策状态}\\
&\to\text{被动观测的秩不足}
\to\text{对偶动作}\\
&\to\text{三条路径的解释}
\to\text{测量隔离的执行}\\
&\to\text{真值隔离的检验}.
\end{aligned}
\]

各包挂在这条链上的位置如下。

\begin{enumerate}
\item 信息赤字由第 1 章定义。它是决策泛函的剩余不确定度，不是传感器清单的长度。
\item 多尺度场由第 4 章写出。已积分的是十六条控制体，不是该章的全部未知量。
\item 决策状态 \(\vect{z}\) 是停机和探针实际使用的泛函。日志里的分量窄于定义。
\item 被动观测的秩由命题 \ref{prop:steady} 说明：平衡点不含内部热阻。
\item 对偶动作在规格里是探针库上的一次比较。日志里是 150 s 至 190 s 的固定窗。
\item 三条路径的符号由结构方程指定。核里路径 C 的结壳不改导热场。
\item 测量隔离的执行在规格里是不读后验的屏障。日志里加热器温度低于削减阈值，所以这条隔离没有在战役中改变指令。
\item 真值隔离的检验要求估计器不读真值。热阻更新使用了由真值温差构成的参照，因此这一环没有严格成立。正文报告方差下降时同时报告这个参照。
\end{enumerate}

文献说明每一环继承了什么。外部数据只收紧淀粉活度这一环的温度上限。演示把秩、路径和投影变成可看的状态，不另写差分。

\paragraph{本节不能推出的结论}
链写全，不能推出八个环节都已经定量成立。上面第 2、5、6、7、8 环都带有“规格”和“日志”两个主语。摘要只用日志那个主语。

\section{风险写在适用区，不写在口号里}

\paragraph{要判定的量}
每一条限制缩小的是哪一句的适用范围。限制不集中成一段无法对应到方程的声明。

\begin{itemize}
\item 两节点命题通过，只允许声明判据方向：平衡点不含内部参数，第二段输入能增加滞后，短窗内冷点增量有上界，安全投影不读后验。主结论等待带 GAB、各向异性和脂肪几何的异构真值。该真值本次没有运行。
\item 等温线没有肉类。肉类活度保持为开放参数。土豆参数只到 \(60^\circ\mathrm{C}\)。
\item 公开温度历程尚未接入。外部预测误差不出现在摘要。
\item 三维切面若由十六条控制体插值，页面上必须能看出它是插值，而不是孔隙尺度的云图。
\item \(\alpha_m=0.85\) 是设计线，对应 \(F\approx 3.79\,\mathrm{min}\)。它不是认证结论。
\item 加热器削减的 \(230^\circ\mathrm{C}\) 与命题函数的 \(210^\circ\mathrm{C}\) 是两段代码。二者都不随后验变形。二者在本次加热器轨迹上都不动作。
\item 热阻更新的参照使用真值温差。方差从 \(0.64\) 乘到 \(0.224\) 不是纯观测滤波器的后验熵。
\end{itemize}

这些限制留在相应章的末尾，并在使用该句的地方再指向一次。它们缩小的是句子的适用范围。未完成的环节不提前写成结果。

\paragraph{本节不能推出的结论}
把限制写进适用区，不能代替把缺的积分器跑完。适用区只防止当前句子越过已经跑完的范围。
